\subsection{Quotation}
\label{Section 2.3.1}

If we can form compound data using symbols, we can have sequences such as

\begin{codeBlock}
>>> ("a", "b", "c", "d")
('a', 'b', 'c', 'd')
>>> (23, 45, 17)
(23, 45, 17)
>>> (("Norah", 12), ("Molly", 9), ("Anna", 7), ("Lauren", 6))
(('Norah', 12), ('Molly', 9), ('Anna', 7), ('Lauren', 6))
\end{codeBlock}

\noindent
There is nothing remarkable about that, and the flatness of the observation is
worth pausing on, because in the original it is not flat at all.  This is the
section where the original introduces a new element into its language, and we
have just done without it.

\begin{buildFromScarcity}
The original needs \newterm{quotation} because in Lisp a program is a list.
Writing \code{(a b c)} in a Lisp program is a procedure call, so there must be
some way to say that you meant the \emph{list} \code{(a b c)} and not the
calling of \code{a}; the mark \code{'} is that way, and \code{'(a b c)} is the
list.  The very uniformity that makes Lisp powerful---that code and data are
the same kind of thing---is what creates the ambiguity that quotation resolves.

Python never merged the two.  \code{("a", "b", "c")} is a tuple and was never
going to be a procedure call, so there is nothing to disambiguate and no mark
to introduce.  What is more, Python already has the device the original reaches
for by analogy: the original observes that natural languages use quotation
marks to indicate that a word is to be treated literally as a string of
characters, and then notes that its own quotation works differently.  Ours does
not.  We write quotation marks around the characters, and get a thing whose
value is those characters.  The new element of this section is, for us, one the
language had from the beginning.

We should be clear about what we are giving up, because it is not nothing.  In
the original, quotation makes every expression available as data, and that is
the foundation of \link{Chapter 4}, where an interpreter for Lisp is written in
Lisp using very little more than the machinery of this section.  A Python
program is not a Python data structure, and no amount of quoting will make it
one.  When we reach that chapter we shall be starting from somewhere quite
different.
\end{buildFromScarcity}

A symbol, for us, is a string: a sequence of characters, written between
quotation marks.\footnote{Which is a looser notion than the original's.  A
Scheme symbol is atomic---it has no parts---while a Python string is a sequence
and will let you ask for its third character or its length.  Nothing in this
chapter does so, and our \code{is\_pair} answers \code{False} for a string, so
a symbol is a leaf wherever we meet one.  But the reader should know that the
indivisibility is a convention we keep rather than a property we have.}

Now we can distinguish between symbols and their values.  The distinction is
the same one the original is making; only the notation is duller.

\begin{codeBlock}
>>> a = 1
>>> b = 2
>>> (a, b)
(1, 2)
>>> ("a", "b")
('a', 'b')
>>> ("a", b)
('a', 2)
\end{codeBlock}

\noindent
The first gives the values of \code{a} and \code{b}, the second the symbols
themselves, the third one of each.  In the original these three are
distinguished by where the quotation mark falls; here they are distinguished by
whether a name is written bare or in quotes, which comes to the same thing.

\begin{codeBlock}
>>> ("a", "b", "c")[0]
'a'
>>> ("a", "b", "c")[1:]
('b', 'c')
\end{codeBlock}

\noindent
\begin{notTaste}
There is a trick worth knowing, and a place where it is the right thing to do.
If we bind a name to the string of its own name,

\begin{codeBlock}
>>> x = "x"
>>> y = "y"
\end{codeBlock}

\noindent
then the name and the symbol coincide, and expressions built from symbols can
be written with bare names:

\begin{codeBlock}
>>> ("*", x, y)
('*', 'x', 'y')
>>> ("+", x, 3)
('+', 'x', 3)
\end{codeBlock}

\noindent
This closes the gap between a name and its value that we were at such pains to
open a moment ago, which is why it must come afterwards: had we begun here, the
three expressions above would all have been the same and there would have been
nothing to see.  Used deliberately, it buys a good deal of readability in
\link{Section 2.3.2}, where the expressions being differentiated are built from
exactly such symbols.

It is a convention and not a mechanism.  Nothing stops anyone writing \code{x =
"y"}, and everything downstream will then be quietly wrong.  Python libraries
that take this seriously do it properly: SymPy's \code{symbols("x")} returns an
object that knows its own name, so the correspondence is established by the
library rather than promised by the programmer.  What we are doing here is the
poor relation of that, and worth the saving only because our symbols never need
to do anything but be themselves.
\end{notTaste}

The original introduces one further primitive for manipulating symbols,
\code{eq?}, which tests whether two symbols are the same.  Ours is \code{==},
and it needs a warning.

\begin{revealTheSurface}
Scheme's \code{eq?} asks whether two symbols are the same object, which invites
the Python reader to reach for \code{is}.  Do not.  \code{is} asks the same
question and gives the right answer often enough to be dangerous:

\begin{codeBlock}
>>> "abracadabra" is "abracadabra"
True
>>> "".join(["abra", "cadabra"]) is "abracadabra"
False
>>> "".join(["abra", "cadabra"]) == "abracadabra"
True
\end{codeBlock}

\noindent
The same characters, and yet the second is false: Python stores one copy of a
string written literally in the source, so two literals are one object, while a
string that was built rather than written is another object with the same
contents.  A symbol comparison written with \code{is} passes every test its
author thinks to try and fails on the first symbol that arrives from somewhere
else.

\code{==} asks whether they have the same characters, which is what the
original means by ``the same'' here.\footnote{The original's own footnote at
this point says that defining sameness as ``the same characters in the same
order'' skirts a deep issue it is not yet ready to address, and promises to
return to it in \link{Chapter 3}.  The distinction between \code{==} and
\code{is} is precisely that issue, arriving early and in a form the reader can
be bitten by.}
\end{revealTheSurface}

Using \code{==}, we can implement a useful procedure called \code{memq}.  This
takes two arguments, a symbol and a sequence.  If the symbol is not contained
in the sequence, then \code{memq} returns false.  Otherwise, it returns the
subsequence beginning with the first occurrence of the symbol:

\begin{codeBlock}
>>> def memq(item, x):
...     if x == ():
...         return False
...     elif item == x[0]:
...         return x
...     else:
...         return memq(item, x[1:])
\end{codeBlock}

\noindent
For example, the value of

\begin{codeBlock}
>>> memq("apple", ("pear", "banana", "prune"))
False
\end{codeBlock}

\noindent
is false, whereas the value of

\begin{codeBlock}
>>> memq("apple", ("x", ("apple", "sauce"), "y", "apple", "pear"))
('apple', 'pear')
\end{codeBlock}

\noindent
is \code{("apple", "pear")}.

Python's \code{in} answers the membership question, so that \code{"apple" in
("pear", "banana")} is \code{False}.  But it answers only that, where
\code{memq} hands back the rest of the sequence from the match onwards.  We
shall want the sequence itself in \link{Section 2.3.3}.

\begin{exerciseGroup}
\phantomsection\label{Exercise 2.53}
\exerciseHeading{Exercise 2.53:} What would the interpreter print
in response to evaluating each of the following expressions?

\begin{codeBlock}
>>> ("a", "b", "c")
>>> (("george",),)
>>> (("x1", "x2"), ("y1", "y2"))[1:]
>>> (("x1", "x2"), ("y1", "y2"))[1:][0]
>>> is_pair(("a", "short", "list")[0])
>>> memq("red", (("red", "shoes"), ("blue", "socks")))
>>> memq("red", ("red", "shoes", "blue", "socks"))
\end{codeBlock}

Two of these are worth more than a moment.  The second is the one the comma
makes or breaks, as \link{Section 2.1.1} warned.  The fifth asks whether a
symbol is a pair, and the answer says something about how we have chosen to
represent symbols.

\phantomsection\label{Exercise 2.54}
\exerciseHeading{Exercise 2.54:} Two sequences are said to be
equal if they contain equal elements arranged in the same order.  For example,

\begin{codeBlock}
>>> ("this", "is", "a", "list") == ("this", "is", "a", "list")
True
>>> ("this", "is", "a", "list") == ("this", ("is", "a"), "list")
False
\end{codeBlock}

\noindent
Python has answered the question already: \code{==} on tuples compares element
by element, and descends into elements that are themselves tuples.  The
original has to build this, defining \code{equal?} recursively in terms of the
\code{eq?} equality of symbols: \code{a} and \code{b} are equal if they are
both symbols and the symbols are the same, or if they are both sequences whose
first elements are equal and whose rests are equal.

Implement it anyway, as \code{is\_equal}, without using \code{==} on anything
but symbols and numbers.  Then say what your procedure does that \code{==} does
not, and what \code{==} does that yours does not.

\phantomsection\label{Exercise 2.55}
\exerciseHeading{Exercise 2.55:} In the original, Eva Lu Ator types
the expression \code{(car ''abracadabra)} and is surprised when the interpreter
prints back \code{quote}.  The explanation is that \code{'x} is an abbreviation
for the list \code{(quote x)}, so \code{''abracadabra} is the list \code{(quote
abracadabra)}, whose first element is the symbol \code{quote}.

There is no corresponding expression in Python, and no surprise to be had.
Explain why not.  Your answer should say what Eva's expression relies on that
ours cannot supply, and it should mention \link{Chapter 4}.
\end{exerciseGroup}
